\clearpage
\section*{三、结果与分析}

\subsection*{1. 油滴电荷量}

重新计算各次测量的电荷量，常量取软件参数。以第 1 次测量为例，$v=l/t=1.5\times10^{-3}/9.99\,\mathrm{m/s}=1.502\times10^{-4}\,\mathrm{m/s}$，未修正半径

\[
r=\sqrt{\frac{9\eta v}{2\rho g}}=\sqrt{\frac{9\times1.83\times10^{-5}\times1.502\times10^{-4}}{2\times981\times9.79}}\ \mathrm{m}=1.135\ \mu\mathrm{m},
\]

修正项 $b/(pr)=6.17\times10^{-6}/(76.0\times1.135\times10^{-6})=0.0716$，代入得 $q_1=4.490\times10^{-18}\,\mathrm{C}$。十次结果如下，均与软件显示值一致。

\begin{center}\small\setlength{\tabcolsep}{4pt}
\begin{tabular}{ccccc}
\toprule
序号 $i$ & 油滴 & $v$/(mm$\cdot$s$^{-1}$) & $r$/$\mu$m & $q_i$/(10$^{-19}$ C) \\
\midrule
1 & 1 & 0.1502 & 1.135 & 44.90 \\
2 & 1 & 0.1459 & 1.119 & 41.54 \\
3 & 2 & 0.3112 & 1.634 & 98.22 \\
4 & 2 & 0.3106 & 1.632 & 96.75 \\
5 & 3 & 0.2555 & 1.480 & 59.04 \\
6 & 3 & 0.2534 & 1.474 & 57.15 \\
7 & 4 & 0.9259 & 2.818 & 224.59 \\
8 & 4 & 0.9146 & 2.800 & 221.59 \\
9 & 5 & 0.5455 & 2.163 & 154.01 \\
10 & 5 & 0.5396 & 2.151 & 151.48 \\
\bottomrule\end{tabular}\end{center}

本次油滴所带电子数为 26～140 个，相邻电荷量之差不再接近 $e$，因此三种方法都取 $e_0=1.602\times10^{-19}\,\mathrm{C}$ 作为参考值，由 $q/e_0$ 取整确定 $n$。

\subsection*{2. 方法一：平均值法}

把电荷量从小到大排列后，同一油滴的两次测量相邻，差值只有 1.5～3.4（单位 10$^{-19}$ C），远小于不同油滴之间 12～68 的差值，因此按油滴分为 5 组。每组取平均 $\bar q$，再求 $n$ 和 $e_n=\bar q/n$。

\begin{center}\small\setlength{\tabcolsep}{4pt}
\begin{tabular}{ccccc}
\toprule
油滴 & $\bar q$/(10$^{-19}$ C) & $\bar q/e_0$ & $n$ & $e_n$/(10$^{-19}$ C) \\
\midrule
1 & 43.22 & 26.98 & 27 & 1.6008 \\
3 & 58.09 & 36.26 & 36 & 1.6137 \\
2 & 97.49 & 60.85 & 61 & 1.5981 \\
5 & 152.74 & 95.35 & 95 & 1.6078 \\
4 & 223.09 & 139.26 & 139 & 1.6050 \\
\bottomrule\end{tabular}\end{center}

五个 $e_n$ 的平均值 $\bar e=1.6051\times10^{-19}\,\mathrm{C}$，样本标准差 $s=0.0061\times10^{-19}\,\mathrm{C}$。

\subsection*{3. 方法二：公约数法}

使用逐次相减法，把 10 个电荷量排序，求相邻差值，再由 $q_i/e_0$ 取整得到 $n_i$，计算 $e_i=q_i/n_i$。

\begin{center}\small\setlength{\tabcolsep}{4pt}
\begin{tabular}{ccccccc}
\toprule
序号 & 原序号 & $q_i$/(10$^{-19}$ C) & $\Delta q=q_{i+1}-q_i$ & $n$ 计算值 & $n$ 取整 & $e_i$/(10$^{-19}$ C) \\
\midrule
1 & 2 & 41.54 & 3.36 & 25.93 & 26 & 1.5978 \\
2 & 1 & 44.90 & 12.25 & 28.03 & 28 & 1.6036 \\
3 & 6 & 57.15 & 1.89 & 35.67 & 36 & 1.5875 \\
4 & 5 & 59.04 & 37.71 & 36.85 & 37 & 1.5956 \\
5 & 4 & 96.75 & 1.48 & 60.39 & 60 & 1.6124 \\
6 & 3 & 98.22 & 53.25 & 61.31 & 61 & 1.6102 \\
7 & 10 & 151.48 & 2.53 & 94.56 & 95 & 1.5945 \\
8 & 9 & 154.01 & 67.58 & 96.14 & 96 & 1.6043 \\
9 & 8 & 221.59 & 3.00 & 138.32 & 138 & 1.6057 \\
10 & 7 & 224.59 & — & 140.19 & 140 & 1.6042 \\
\bottomrule\end{tabular}\end{center}

十个 $e_i$ 的平均值 $\bar e=1.6016\times10^{-19}\,\mathrm{C}$，$s=0.0076\times10^{-19}\,\mathrm{C}$。

\subsection*{4. 方法三：多重差值法}

先对排序后的电荷量求一阶差值 $\Delta q_1$，再对 $\Delta q_1$ 求二阶差值 $\Delta q_2=|\Delta q_{1,j+1}-\Delta q_{1,j}|$。

\begin{center}\small\setlength{\tabcolsep}{4pt}
\begin{tabular}{cccccccccc}
\toprule
序号 $j$ & 1 & 2 & 3 & 4 & 5 & 6 & 7 & 8 & 9 \\
\midrule
$\Delta q_1$/(10$^{-19}$ C) & 3.36 & 12.25 & 1.89 & 37.71 & 1.48 & 53.25 & 2.53 & 67.58 & 3.00 \\
$k$ & 2 & 8 & 1 & 24 & 1 & 33 & 2 & 42 & 2 \\
$e$/(10$^{-19}$ C) & 1.678 & 1.531 & 1.886 & 1.571 & 1.478 & 1.614 & 1.265 & 1.609 & 1.499 \\
\bottomrule\end{tabular}\end{center}

\begin{center}\small\setlength{\tabcolsep}{4pt}
\begin{tabular}{ccccccccc}
\toprule
序号 $j$ & 1 & 2 & 3 & 4 & 5 & 6 & 7 & 8 \\
\midrule
$\Delta q_2$/(10$^{-19}$ C) & 8.89 & 10.36 & 35.82 & 36.23 & 51.78 & 50.72 & 65.05 & 64.58 \\
$k$ & 6 & 6 & 22 & 23 & 32 & 32 & 41 & 40 \\
$e$/(10$^{-19}$ C) & 1.482 & 1.727 & 1.628 & 1.575 & 1.618 & 1.585 & 1.587 & 1.615 \\
\bottomrule\end{tabular}\end{center}

17 个值的平均值 $\bar e=1.585\times10^{-19}\,\mathrm{C}$，$s=0.127\times10^{-19}\,\mathrm{C}$。可以看到，$k$ 只有 1～2 的差值算出的 $e$ 分散很大（1.265～1.886），因为这些差值与测量误差同一量级。

\subsection*{5. 结果与不确定度}

对多次得到的 $e$ 取平均，A 类标准不确定度为 $u_A=s/\sqrt N$。电压和计时均由仪器数字显示，现场没有给出它们的允许误差限，因此不评定 B 类分量。

\begin{center}\small\setlength{\tabcolsep}{4pt}
\begin{tabular}{ccccccc}
\toprule
方法 & $N$ & $\bar e$/(10$^{-19}$ C) & $s$/(10$^{-19}$ C) & $u_A$/(10$^{-19}$ C) & 测量结果/(10$^{-19}$ C) & 相对误差 $E$ \\
\midrule
平均值法 & 5 & 1.6051 & 0.0061 & 0.0027 & $1.6051\pm0.0027$ & +0.18\% \\
公约数法 & 10 & 1.6016 & 0.0076 & 0.0024 & $1.6016\pm0.0024$ & −0.04\% \\
多重差值法 & 17 & 1.585 & 0.127 & 0.031 & $1.59\pm0.03$ & −1.05\% \\
\bottomrule\end{tabular}\end{center}

其中 $E=(\bar e-e_{\text{公认}})/e_{\text{公认}}\times100\%$，$e_{\text{公认}}=1.602177\times10^{-19}\,\mathrm{C}$。三种方法的结果分别为

\[
e_1=(1.6051\pm0.0027)\times10^{-19}\ \mathrm{C},
\]

\[
e_2=(1.6016\pm0.0024)\times10^{-19}\ \mathrm{C},
\]

\[
e_3=(1.59\pm0.03)\times10^{-19}\ \mathrm{C}.
\]

平均值法和公约数法与公认值的偏差分别约为 1.1 倍和 0.25 倍不确定度，测量结果与公认值吻合较好；多重差值法中小差值的相对误差较大，不确定度最大。

\subsection*{6. 误差分析}

\begin{enumerate}
\item \textbf{油滴带电量较多。} 本次 $n=26\sim140$，相邻电荷量的差值较大，需借助 $e_0$ 确定 $n$；同时小差值与测量误差同一量级，使多重差值法的结果较为分散。
\item \textbf{计时误差。} 由 $q\propto t^{-3/2}$，下落时间的相对误差会放大 1.5 倍传到 $q$。第 4 粒油滴下落时间只有约 1.6 s，按键计时的反应误差影响较大。
\item \textbf{平衡电压判断。} 平衡电压以整数伏读取，油滴是否完全静止靠肉眼判断，因为布朗运动，油滴难免小范围晃动，难以完美达到静止平衡。
\item \textbf{油滴变化。} 本次油滴半径约 1.1～2.8 $\mu$m，较小的油滴受空气分子碰撞会上下抖动，测量过程中油滴可能少量挥发。

\end{enumerate}
